\begin{abstract}
This study develops and retrospectively evaluates a probabilistic simulation
framework for the FIFA World Cup 2026. The repository combines official squad
records, EAFC-style player ratings, public player statistics and national-team
strength priors to simulate complete tournaments. Three primary approaches are
compared: a Poisson baseline, a heuristic event simulator, and the same event
simulator augmented with a trained Random Forest. A legacy Random Forest run
is retained as a reference. Each saved forecast contains 10,000 tournament
realisations. Predictions are compared with FIFA's published outcomes using
stage-wise Brier scores, champion log loss and top-$k$ qualifier overlap.
Spain, the eventual champion, received title probabilities of 10.26\%, 18.18\%
and 17.11\% in the primary models, although none ranked Spain first. Heuristic
v2 achieved the lowest mean six-stage Brier score, 0.0840, compared with 0.0872
for the baseline and 0.0867 for Random Forest v2. Both v2 approaches ranked
Kylian Mbapp\'{e} first by expected goals. The analysis also identifies substantial
errors in Germany's and Brazil's projected advancement and in the prospects
of Norway and Morocco. Interpretation is limited by heterogeneous player-data
coverage, synthetic model training, and the absence of independently audited
pre-tournament publication timestamps. The results support the value of a
transparent simulation pipeline while motivating historical backtesting,
probability calibration and better source-version tracking.
\end{abstract}
\noindent\textbf{Keywords:} Monte Carlo simulation; football forecasting;
Poisson model; Random Forest; Brier score; data provenance.
\clearpage
\renewcommand{\contentsname}{Contents}
\tableofcontents
\clearpage
\listoffigures
\listoftables
\clearpage

\section{Introduction and Project Objectives}
\label{sec:introduction}
Tournament forecasting involves several sources of uncertainty: which players
are available, how team strength translates into match outcomes, and how the
draw and knockout path shape progression. A team can have the greatest chance
of winning while still being more likely to lose the tournament than to win it.
The project's central output is therefore a distribution of possible outcomes.

The repository began with an interpretable baseline and subsequently added
player-performance information, team-rating priors and a more detailed match
engine. Keeping the baseline makes it possible to ask whether additional
complexity improves the forecasts, and where it fails to do so.
\subsection{Research Questions}
\begin{enumerate}
\item How much forecast probability did each model assign to the eventual champion?
\item How accurately did the models allocate advancement probabilities across all 48 teams?
\item Did the event simulator and learned outcome layer improve on the simpler baseline?
\item Which unexpected results reveal weaknesses in team strength, tournament routing or data coverage?
\item Which conclusions are supported by saved evidence, and which require further validation?
\end{enumerate}
\subsection{Scope and Contributions}
The project covers a 48-team competition with 12 groups of four, followed by
a 32-team knockout phase. The top two teams in each group and eight best
third-place teams progress. Outputs include title and stage probabilities,
player goal and assist summaries, match-event totals, model explanations and
data-quality reports.

This report brings together the data pipeline, model definitions and evaluation
of the saved forecasts. The comparison uses stored results without retraining
on the observed tournament. The current source code documents the implemented
mechanisms, but its presence alone does not establish an immutable code version
for every June run. Frozen JSON metadata are the primary evidence for each
run's inputs, seed and model mode.
\subsection{Observed Tournament Outcome}
Spain defeated Argentina 1--0 after extra time in the final. England finished
third and France fourth. Mbapp\'{e} won the Golden Boot with 10 goals
\cite{outcome-final,outcome-standings,outcome-awards}. These observations supply
evaluation targets; they are not used to update the frozen forecasts.
\clearpage

\section{Datasets, Sources and Preparation}
\label{sec:data}
The project uses several datasets with different coverage, units and time
granularities. A player appearing in the combined file does not imply that
every player has a complete, comparable recent-season record.
\subsection{Official Squads}
The working roster contains 48 teams and 1,248 players. Its metadata identify
FIFA's English squad-list PDF, Version 1, dated 3 June 2026 at 11:30 UTC, as the
source \cite{data-fifa-squads}. The script
\path{rebuild_squads_from_fifa_pdf.py} extracts squad records; the active output
is \path{guardian_world_cup_2026_player_guide.json}. Despite its inherited
filename, the current file identifies FIFA as its squad source. The Guardian
guide is separately listed in tournament-structure provenance \cite{data-guardian}.

The source URL is a live resource: when checked for this report, the online PDF
showed a later July timestamp. That page cannot, by itself, prove the exact
contents of the locally recorded June snapshot. Archive the original file and
its checksum when reproducing or extending the study.
\subsection{EAFC Player Ratings}
The local \path{eafc26_players.csv} contains 18,405 records. Its update-date
column is 19 September 2025 throughout the file. It supplies player names,
nationality, positions, overall ratings and additional attributes. The saved
baseline reports 980 matched players and 268 fallback ratings out of 1,248:
approximately 78.5\% matched coverage and 21.5\% fallback coverage.

The export includes SoFIFA-hosted image URLs and player-path fields. These are
provenance clues, not proof of the file's original distributor. The repository
does not preserve a verified download URL, dataset release identifier or licence
for this exact CSV. Accordingly, this report cites the local artifact and
records that upstream provenance gap explicitly \cite{data-eafc}.
\subsection{Public Player-Performance Data}
The v2 input is \path{player_performance_data_statbunker.json}. Its metadata
record 19,923 source rows and 1,248 output players. Source rows can repeat a
player across competitions, seasons and table types; they are not 19,923 unique
players. StatBunker supplies the main public competition-table input, including
appearances, goals, assists, cards and goalkeeper clean sheets \cite{data-statbunker}.

Transfermarkt-derived CSVs from \path{salimt/football-datasets} supplement player
identity and performance information \cite{data-transfermarkt}. English Wikipedia
biographies and manual public-web research fill some remaining gaps
\cite{data-wikipedia}. Two players have SofaScore-derived aggregate statistics
recorded in the final metadata \cite{data-sofascore}. The final source-summary
counts are shown in Table~\ref{tab:coverage}.
\clearpage
@@SOURCE_TABLE@@
\subsection{Coverage and Missingness}
StatBunker's default pages do not supply complete minutes, expected goals,
expected assists, progressive passing, tackles, saves or match ratings. Wikipedia
infoboxes often contain career totals rather than recent-season splits.
Transfermarkt supplements do not make all records equivalent in detail or
freshness. The metadata's ``complete'' flag describes populated player records;
two partial-context rows remain documented.

The research coverage report illustrates the unevenness. Panama has 18 players
below the stored 0.55 confidence threshold and 24 with zero recorded minutes;
Senegal has 14 and 25 respectively. A zero in these source-derived fields can
reflect missing evidence rather than a player having played no football.
The stored confidence measure is an engineering heuristic, not a calibrated
probability that the record is correct.

An earlier squad-quality report recorded 17 squad-shape issues and a
\texttt{needs\_refresh} status. It predates later enrichment, so it should be
treated as a historical diagnostic requiring a fresh check, not proof that all
17 issues remain in the current roster.
\subsection{National-Team Rating Priors}
The primary v2 forecasts used FIFA ranking positions transformed to an
Elo-like scale \cite{data-fifa-ranking}:
\begin{equation}
R_{\mathrm{prior}}=2170-130\ln(r),
\end{equation}
where $r$ is ranking position. The preserved fallback file records a ranking
date of 19 November 2025. This is a model-defined transformation of ranking
positions, not FIFA ranking points and not a historical World Football Elo series.

A subsequent import brought in 177 World Football Elo team records with an
as-of date of 11 June 2026 \cite{data-elo}. The import report shows coverage of all
required teams. These newer inputs are present locally, but the frozen v2 runs
evaluated here still declare the older rank-derived prior. The current input
file and the input version used by a saved result must therefore be distinguished.
\clearpage
\subsection{Tournament and Venue Context}
Groups and tournament rules are stored in
\path{world_cup_2026_simulation_structure.json}, whose metadata cite FIFA draw
information and the Guardian player guide \cite{data-draw,data-guardian}.
\path{world_cup_2026_venue_context.json} provides group-region and stage-level
context. Heat, humidity and travel-load values are modelling priors; they are
not measured match-day weather observations or a complete travel itinerary.
\subsection{Identity Resolution and Enrichment}
The preparation pipeline has four steps:
\begin{enumerate}
\item Extract squad identities, broad positions and clubs from the official list.
\item Normalize names and compare candidates using nationality, name similarity,
exact aliases and club information. Review overrides are stored separately.
\item Attach ratings and detailed positions. For unmatched EAFC players, use
role-based defaults: 68 for goalkeepers/defenders and 69 for midfielders/forwards.
\item Merge public statistics with source identity, collection status and confidence,
then construct team profiles from player records.
\end{enumerate}
Name normalization supports matching across sources; presentation should retain
the player's proper spelling. The report therefore uses Kylian Mbapp\'{e}.
Some archived records retain unaccented matching keys; those frozen identifiers
are preserved to keep the evaluation reproducible.

\path{performance_data_config.json} specifies intended collection windows,
competition weights and stage multipliers. These settings describe collection
and modelling intentions. A listed field or provider is not evidence that it
was available for every player or consumed by every saved simulation.
\subsection{Attempted and Future Data Sources}
SofaScore access was explored through direct and browser-assisted collectors;
the final source summary records only two primary SofaScore rows. Sportmonks,
API-Football/API-Sports and FootyStats have provider-ready collection paths,
but the saved input summary does not identify them as primary suppliers.

The project also attempted to fetch the international-results dataset maintained
by \path{martj42/international_results} \cite{data-historical}. The June manifest
records a failed download and no normalized results file. The stored Random
Forest training metadata consequently report zero historical rows. This
dataset is a planned input for supervised backtesting, not training evidence
for the forecasts evaluated here.
\subsection{Dataset Register}
The accompanying \path{dataset_sources.csv} records each source, the corresponding
local artifact, how it was used, and its provenance limitations. Provider URLs
are distinguished from verified URLs for the exact snapshot. The register
should be updated whenever inputs change.
\clearpage

\section{Models and Simulation Architecture}
\label{sec:models}
\subsection{Monte Carlo Tournament Simulation}
Monte Carlo simulation repeatedly draws match outcomes and plays through the
competition. For event $A$, its estimated probability after $N$ tournaments is
\begin{equation}
\widehat{P}(A)=\frac{1}{N}\sum_{i=1}^{N}\mathbf{1}\{A\text{ occurs in run }i\}.
\end{equation}
Here $N=10{,}000$ for each compared artifact. This is a count of simulated
tournaments, not machine-learning training epochs. Under independent runs,
the conditional Monte Carlo standard error is approximately
$\sqrt{\widehat p(1-\widehat p)/N}$. For a 20\% title probability it is about
0.4 percentage points. Running more simulations reduces this sampling noise;
it does not repair biased ratings or inaccurate football assumptions.

\subsection{Model 1: EAFC--Poisson Baseline}
The baseline in \path{simulate_world_cup.py} builds a position-aware starting
XI, a top-18 group and a full-squad mean. Team strength is
\begin{equation}
S=0.62\overline R_{\mathrm{XI}}+0.28\overline R_{18}
  +0.10\overline R_{\mathrm{squad}}.
\end{equation}
For strength difference $\Delta=S_A-S_B$, expected goals are
\begin{align}
\lambda_A&=\operatorname{clip}(1.28+0.045\Delta,\,0.25,\,3.4),\\
\lambda_B&=\operatorname{clip}(1.28-0.045\Delta,\,0.25,\,3.4).
\end{align}
The match score is sampled using Poisson goal counts. The baseline is useful
because its rating-to-score relationship is simple to inspect. Its limitations
include static strength, role-based fallback ratings and limited representation
of changing match state. Its stored metadata describe knockout seeding based on
group-stage performance, whereas v2 uses constrained bracket routing; differences
between model results therefore reflect more than the learning algorithm alone.

\subsection{Shared v2 Event Engine}
The v2 implementation adds position-sensitive player-performance adjustments,
lineup selection and separate attack, midfield, defence, goalkeeper and
possession profiles. Match context includes team priors, host advantage,
rest, accumulated fatigue, lineup replacements and venue/travel assumptions.
Regulation time is compressed into nine simulated ticks and extra time into
three. Possession, fouls, cards, removals, substitutions, set pieces and game
state can affect scoring; tied knockout matches proceed to extra time and
penalties. The engine carries player state through the tournament and supports
CPU-parallel tournament batches.
\clearpage
\subsection{Model 2: Heuristic Event Simulator}
The heuristic mode uses 96 deterministic, seeded rule-based votes to adjust
attack share, tempo, discipline and possession. The thresholds are hand-built
mechanics, not trees learned from observed match outcomes. The aggregate
modifiers feed the expected-goal calculation before the event engine produces
the match result.

This version asks whether additional team and match context improves on the
baseline without fitting a supervised outcome classifier. It is more expressive,
but the additional coefficients and constraints also introduce assumptions
requiring validation. Both v2 modes retain this heuristic component.
\subsection{Model 3: Random Forest-Augmented Event Simulator}
The learned component is a scikit-learn Random Forest classifier with outcomes
$A$ win, draw and $B$ win. The inspected implementation uses 96 trees, maximum
depth 8, at least 6 samples per leaf, balanced-subsample class weighting and a
fixed seed. Its 21 inputs cover strength and unit differences, discipline and
data confidence, host/rest/fatigue context, venue/travel conditions and tournament
stage information. A model feature need not represent directly measured evidence:
several inputs are derived priors.

The forest modifies the event simulator rather than replacing it. Let
$s_{\mathrm{base}}$ be the engine's attacking share and $p_A,p_D,p_B$ the forest's
outcome probabilities. The pre-clipping blend is
\begin{equation}
s_A=0.72s_{\mathrm{base}}+
0.28\frac{p_A}{\max(0.01,p_A+p_B)}.
\end{equation}
The draw probability also adjusts total expected goals. Consequently, the final
tournament probabilities reflect the combined rating model, heuristic rules,
learned adjustment and stochastic event engine.
\subsection{Synthetic Training and Its Implications}
The stored forest uses 6,768 synthetic rows and no historical match rows. This
count follows from $48\times47$ ordered team pairings across three contexts.
The generator perturbs features and samples outcome labels from a hand-defined
probability function involving team strength, unit matchups and context.

The model therefore learns the behaviour of the synthetic generator. Good
performance on held-out synthetic rows does not establish real-match accuracy.
The saved RF run's 0.7301 accuracy is measured on its training data and must not
be presented as a 73.01\% World Cup prediction success rate. The code supports
historical training when enough usable labels exist; that path was not active
for these artifacts.
\clearpage
\subsection{Legacy Reference and Experimental Controls}
The legacy RF run uses squad-proxy Elo rather than the national-team rank prior.
It is retained to expose the effect of earlier assumptions, but excluded from
the primary model-average comparison. Different seeds, bracket mechanics,
priors and context versions mean the stored runs are not a controlled ablation
of one factor at a time.
@@RUN_TABLE@@
\subsection{Additional Research Models}
The research pipeline also benchmarks L2 logistic regression, sigmoid-calibrated
logistic regression, Extra Trees, histogram gradient boosting and an ensemble
weighted by inverse validation log loss. Logistic regression provides a linear
reference; calibration adjusts probability outputs; the tree ensembles explore
nonlinear feature interactions. These are diagnostic experiments, not five
additional frozen tournament forecasts in this retrospective comparison.

The stored research split contains 5,076 training rows and 1,692 test rows from
the synthetic frame. Calibrated logistic regression has log loss 0.75862;
Random Forest has 0.79449. These numbers use synthetic targets and are not
comparable to champion log loss on the actual tournament.

The ensemble chooses its weights using the same held-out results on which the
ensemble is then scored. Its reported score is consequently not an independent
final test. Future work should separate fitting, calibration, model selection
and final chronological evaluation. Random splitting of related synthetic
team pairings is also weaker evidence than held-out tournaments or dates.
\subsection{Feature Importance and Explainability}
The earlier feature study uses tree importance, permutation importance,
feature-group ablation, partial dependence, individual conditional expectation
and local feature occlusion. Squad-rating difference, possession difference and
attack-versus-defence are leading signals in the stored permutation results.
Unit-matchup ablation yields the largest reported validation log-loss increase,
approximately 0.0228.

These results explain sensitivity within the synthetic modelling framework;
they do not establish causal determinants of real football outcomes. The
tournament-level RF ablation uses only 30 runs per scenario and leaves the
hand-built expected-goal mechanics active. Those exploratory results are not
equivalent to full-engine ablation or a 10,000-run controlled comparison.
\clearpage
